% Pista alg-23j2-q2 · Àlgebra
% Procedència: PAU Catalunya 2023 · Convocatòria de juny · Sèrie 5 · Qüestió 2
\documentclass[11pt,a4paper]{article}
\usepackage{pista}
\pistainfo{Catalunya 2023}{Convocatòria de juny}{Sèrie 5}

\begin{document}

\section*{\textcolor{blauPista}{Qüestió 2}}

\noindent Considereu el sistema d'equacions lineals
\[
\left.\begin{array}{r}
x - y + kz = -1 \\
x + ky + z = 3 \\
2x + (k - 1)y + 2z = 2
\end{array}\right\},
\]
en què $k$ és un paràmetre real.

\subsection*{\textit{a)} \normalfont Discutiu el sistema en funció del valor de $k$. \hspace*{\fill}\textnormal{[1,5~punts]}}

\pista{Esglaona la matriu ampliada, aplicant el mètode de Gauss. En primer lloc, has d'aconseguir dos zeros a la primera columna:
\[
F_{2} \to F_{2} - F_{1}, \qquad F_{3} \to F_{3} - 2F_{1}
\]
Fixa't que les noves files 2 i 3 comencen igual. Per això, amb $F_{3} \to F_{3} - F_{2}$ eliminaràs la $y$ de la tercera fila.

\medskip
Ara mira la diagonal: per a quins valors de $k$ hi surt algun zero? Són els casos especials. Estudia'ls d'un en un, i la resta de valors es tracten tots en un sol cas, perquè el sistema és determinat i té solució única.

\medskip
\textbf{Atenció:} en algun cas especial, potser caldrà un pas més de Gauss abans de comptar rangs.}

\subsection*{\textit{b)} \normalfont Resoleu el sistema per a $k = 0$ i per a $k = 1$. \hspace*{\fill}\textnormal{[1~punt]}}

\pista{Aprofita la matriu esglaonada de l'apartat \textit{a)}: substitueix-hi el valor de $k$ i resol les equacions, començant per la 3a equació, després la 2a equació i finalment la 1a equació.

\medskip
Si amb algun dels dos valors et queda una fila tota de zeros, tindràs més incògnites que equacions. Llavors, pren-ne una com a paràmetre, per exemple $z = t$, i expressa les altres en funció de $t$.}

\end{document}
